% Zdroj: PDF priklady-jaro-2026.pdf, fyzická str. 22. Značka: specializace UI-SU, UI-ZPJ. Zařazeno: S3.
\section{Kombinace modelů}
\szzotazka{jaro 2026}{29}{Kombinace modelů}

\noindent\textbf{Zadání.}\par
\begin{enumerate}
  \item Stručně popište základní myšlenky technik bagging a boosting. Uveďte příklad konkrétních modelů, které tyto techniky
    používají.
  \item Uvažujme tři různé modely $M_1$, $M_2$ a $M_3$ pro binární klasifikaci. Přesnost (accuracy) modelu $M_1$ je $0{,}70$, přesnost
    modelu $M_2$ je $0{,}75$ a přesnost modelu $M_3$ je $0{,}80$. Z těchto modelů vytvoříme ensemble tak, že modely necháme hlasovat
    (vzorek se předloží všem modelům a jako finální odpověď se použije ta, která je častější). Za jakých podmínek bude
    přesnost výsledného ensemble lepší, než přesnost nejlepšího modelu? Jaká je nejlepší dosažitelná přesnost ensemblu? Za
    jakých podmínek bude naopak přesnost ensemblu horší než přesnost nejlepšího modelu? \emph{Nápověda: Uvažujte předpovědi
    každého z modelů na jednotlivých vzorcích dat.}
\end{enumerate}

\medskip
\noindent\textbf{Náčrt řešení.}\par
\begin{enumerate}
  \item Bagging trénuje modely na různých náhodných podmnožinách trénovací množiny (bootstrap). Výstupy modelů se
    následně agregují (např. hlasováním) pro finální predikci ensemblu. Při boostingu se trénuje posloupnost modelů tak,
    že instance špatně předpovězené předchozím modelem mají větší váhu. Příkladem baggingu může být random forest,
    příkladem boostingu AdaBoost nebo XGBoost.
  \item Pokud každý z modelů chybně předpovídá jiná data (množiny chybně předpovězených instancí jsou navzájem dis\-junktní),
    výsledný ensemble bude mít accuracy $100\,\%$. Pokud by naopak množina instancí chybně klasifikovaná
    modelem $M_2$ byla podmnožinou množiny chybně klasifikovanou $M_1$, bude accuracy přinejlepším $0{,}75$. Pokud by navíc
    chybně klasifikované vzorky modelu $M_3$ pokryly zbytek chybně klasifikovaných modelem $M_1$, bude accuracy přinejlepším
    $0{,}7$.
\end{enumerate}

\medskip
\noindent\textit{Doplňující vysvětlení (nad rámec oficiálního náčrtu):} Označme $E_i$ množinu vzorků, na nichž chybuje $M_i$ ($|E_1|=30\,\%$, $|E_2|=25\,\%$, $|E_3|=20\,\%$). Ensemble se mýlí právě na vzorcích, které leží aspoň ve dvou množinách $E_i$. Je lepší než nejlepší model, právě když takových vzorků je méně než $20\,\%$, tj.\ když se chybové množiny překrývají málo (chyby modelů jsou málo korelované). Při nezávislých chybách by vyšla accuracy $1-(0{,}14+0{,}015)=0{,}845$. Horší než nejlepší model je ensemble při velkém překryvu; nejhorší možná accuracy je $1-\frac{0{,}30+0{,}25+0{,}20}{2}=0{,}625$ (každý chybný vzorek ensemblu spotřebuje aspoň dvě chyby modelů), např.\ při $|E_1\cap E_2|=17{,}5\,\%$, $|E_1\cap E_3|=12{,}5\,\%$, $|E_2\cap E_3|=7{,}5\,\%$ a~prázdném trojném průniku.
